% ===== 第 1 章：问题重述 / 问题背景 / 假设与判据口径 / 符号说明 =====
% 对应 blueprints/paper_blueprint.md 章节树 01_intro_assump。
% 内容：问题重述要点、背景、模型假设（编号+理由）、判据口径溯源、符号说明表、数据/题面参数表。
% 主节编号方式：使用 main_template.tex 中的 \ctexset 自动编号（方式一）则直接
%   \section{问题重述}。若用手工编号则 \section*{一、问题重述}。
\section{问题重述}

\subsection{问题背景}
[在此写入：题面背景、问题背景要点、相关物理/机理图像（先直观后公式）。]

\subsection{需要解决的问题}
\begin{itemize}
  \item \textbf{问题一}：[问题一描述]。
  \item \textbf{问题二}：[问题二描述]。
  % ... 按题面补全，列出所有问题
\end{itemize}

\subsection{模型假设}
\begin{enumerate}
  \item[1.] \textbf{[题面明示/合理简化]} [假设 + 理由]。
  \item[2.] ...  % 每条假设编号 + 标来源（题面明示/合理简化）
\end{enumerate}

\subsection{判据口径}  % 若题目含条件/判定语义，保留；纯评价预测可删
[在此写入：关键判据的溯源、多解性讨论、最终采纳口径。若无可删本节。]

\subsection{符号说明}
\begin{table}[H]
\centering
\caption{主要符号说明}\label{tab:symbols}
\begin{tabular}{cl}
\toprule
符号 & 含义 \\
\midrule
% 按需填符号，例如 $x$ & 说明 \\
\bottomrule
\end{tabular}
\end{table}
